\section{Marco jurídico: a LGPD e o contraponto europeu}

% Objetivo da seção: reconstruir o direito à eliminação na LGPD como eixo
% principal, identificar o "atacante" que a lei pressupõe ao definir dado
% anonimizado, e só então usar o direito europeu como contraponto.

Esta seção deve partir da LGPD. Primeiro, o art.~12 e o teste de "meios
técnicos razoáveis e disponíveis na ocasião do tratamento", que fixa o
adversário contra o qual a anonimização é medida. Depois, o art.~18 e seu
inciso~IV, que coloca anonimização, bloqueio e eliminação como remédios
alternativos à disposição do titular. Por fim, o art.~16, que autoriza a
conservação de dados após o término do tratamento quando a eliminação for
tecnicamente inviável, cláusula sem equivalente direto no RGPD. Deve incluir a
posição da ANPD (guia de anonimização; eventual regulamentação dos
arts.~18 a 20) e a doutrina brasileira sobre dado anonimizado e sobre os
limites do art.~16.

Em seguida, como contraponto, o direito europeu: art.~17 do RGPD,
Considerando~26 (critério dos "meios razoavelmente suscetíveis de serem
utilizados"), Parecer~28/2024 do EDPB sobre quando um modelo de IA pode ser
considerado anônimo, e o acórdão do TJUE em EDPS~v~SRB (C-413/23~P) sobre a
relatividade do conceito de dado pessoal. O objetivo do contraponto é
mostrar que o adversário relativo já é reconhecido na Europa e que a LGPD,
ao prever o art.~16 e o art.~18, IV, está em posição melhor para acolher o
padrão proposto.

\todo{Texto literal dos arts.~12, 16 e 18 em notas ou apêndice.}
\todo{ANPD: guia de anonimização e regulamentação dos arts.~18--20 (ver notes/perguntas-em-aberto.md).}
\todo{Doutrina brasileira: ver lista de candidatos em notes/perguntas-em-aberto.md.}
\todo{Doutrina europeia: Juliussen et al.; Colonna e Oechtering.}
